\section{Qiushi Engine and the development of the proof}\label{sec:development}

Qiushi Engine developed the structural proof through sustained
autonomous research on exact matrix multiplication. Its research
planning, construction, computation, and critical review operate with
persistent mathematical memory, as described in its earlier work on
an optical experimental platform~\cite{qiushi-optics}. Here the system
began with the search for a shorter algorithm. The eventual lower-bound
proof grew from a different question exposed by that search: which
constraints on candidate factors express compatibility with the full
multiplication tensor?

The public research trajectory~\cite{qiushi-report} brings together
the relevant notes, programs, results, and corrections. It is organized
by mathematical question rather than as a transcript or a strictly
linear chronology. The following three developments explain how the
representations changed and what was retained from each.

\subsection{From numerical dependence to an exact deletion criterion}
A general $22$-term ansatz contains $22(9+9+9)=594$ coefficients
subject to $729$ cubic tensor-coordinate identities. The early search
also considered real-coefficient deformations of known $23$-term
schemes. The aim was to reach a configuration in which one product
could be eliminated. A nearly dependent first--second factor pairing
initially appeared to offer such a route.

The difficulty can be stated exactly. For a numerical decomposition
$\mathcal D=\sum_{t=0}^{22}u_t\otimes v_t\otimes w_t$ over
$\mathbb R$, normalize the nonzero pairing columns by setting
\[
 d_t=\frac{u_t\otimes v_t}{\|u_t\otimes v_t\|_2},\qquad
 \widetilde w_t=\|u_t\otimes v_t\|_2 w_t,\qquad
 D=[d_0\ \cdots\ d_{22}].
\]
Let $\lambda$ be a unit right singular vector for
$\sigma_{\min}(D)$, and choose $k$ with $\lambda_k\ne0$.
If $D\lambda$ were zero, the $k$-th pairing column could be
eliminated by replacing each remaining output vector with
$\widetilde w_t-(\lambda_t/\lambda_k)\widetilde w_k$.
For an approximate dependence the same operation gives a $22$-term
tensor $\mathcal D_{22}$, with exact error relative to the current
tensor
\begin{equation}\label{eq:deletion-diagnostic}
 \|\mathcal D-\mathcal D_{22}\|_F
   =\frac{\|D\lambda\|_2\,\|\widetilde w_k\|_2}{|\lambda_k|}
   =\frac{\sigma_{\min}(D)\,\|\widetilde w_k\|_2}{|\lambda_k|}.
\end{equation}
Thus a small pairing singular value alone is insufficient. The output
factor may grow at the compensating rate.

To compare the result with the target tensor, write
$E=\|\mathcal D-\mathcal D_{22}\|_F$ and
$\varepsilon=\|\T-\mathcal D\|_F$, using the real tensor with the
coefficients in~\eqref{eq:tensor}. The triangle inequalities give
\begin{equation}\label{eq:deletion-target}
 |E-\varepsilon|\leq\|\T-\mathcal D_{22}\|_F
                  \leq E+\varepsilon.
\end{equation}
Table~\ref{tab:continuation} reports both quantities for a continuation
from the integer $23$-term scheme \code{serendipitous\_8d34} in
Perminov's FastMatrixMultiplication collection~\cite{perminov-code}.
The output vectors were transposed from the source's column-major
convention to~\eqref{eq:tensor}; the source version and conversion
are specified in Appendix~\ref{app:trajectory-sources}.
The parameter $t$ labels saved continuation points, not elapsed time.
The same column, labelled $10$ in the data, is eliminated at each
point. Across these samples the pairing singular value falls by two
orders of magnitude, but output growth keeps the deletion error of
order one. The target residual remains below $4\times10^{-12}$.
Thus~\eqref{eq:deletion-target} shows that the deleted representation
is still about one unit from the target at each of these points.

\begin{table}[htbp]
\centering
\caption{Pairing dependence and compensating output growth in a
recorded real-coefficient continuation. Here $E$ is the deletion error
in~\eqref{eq:deletion-diagnostic}, and $\varepsilon$ is the target
residual before deletion. Values are rounded from the archived data.}
\label{tab:continuation}
\small
\begin{tabular}{rrrrrr}
\toprule
$t$ & $\sigma_{\min}(D)$ & $|\lambda_{10}|$ & $\|\widetilde w_{10}\|_2$ & $E$ & $\varepsilon$\\
\midrule
10 & $5.265\times10^{-3}$ & $0.7084$ & $1.400\times10^2$ & $1.040$ & $1.997\times10^{-14}$\\
20 & $7.051\times10^{-4}$ & $0.7076$ & $1.064\times10^3$ & $1.061$ & $1.902\times10^{-13}$\\
60 & $4.745\times10^{-5}$ & $0.7075$ & $1.779\times10^4$ & $1.193$ & $3.645\times10^{-12}$\\
\bottomrule
\end{tabular}
\end{table}

Qiushi's correction separated projective proximity from the tensor
contribution it was meant to remove. The retained lesson was an exact
one: a change to the first two factor lists must be assessed together
with the output coefficients. The archive includes the continuation
analysis, its corrected interpretation, and the data used in
Table~\ref{tab:continuation}; the source links are collected in
Appendix~\ref{app:trajectory-sources}.

\subsection{Quotient cores and the completion problem}
The finite-field search made a coordinate quotient a natural target.
Taking $W=\langle E_{00}\rangle$ in the first factor gives an
$8\times9\times9$ core. This is the $E_{11}$ core of the research
report, whose historical name uses one-based indices; the present
paper uses indices $0,1,2$. A $19$-term decomposition of this core,
embedded in the coordinate complement, would need only the three
deleted terms
\begin{equation}\label{eq:core-restoration}
 \sum_{k=0}^2 E_{00}\otimes E_{0k}\otimes E_{0k}
\end{equation}
to recover $\T$. It would therefore give a $22$-product algorithm.
The smaller object retained a direct route to the original problem.

Occupation constraints organized candidate first factors into finite
configurations. Their feasibility, however, did not solve completion.
For fixed first factors $a_t\in V/W$, the remaining obligation is
\begin{equation}\label{eq:core-completion}
 \sum_t a_t B_t[j,k]C_t[i,l]
   =\delta_{kl}\,q_W(E_{ij})
 \qquad(0\leq i,j,k,l<3).
\end{equation}
Each pair $B_t,C_t$ must satisfy all these equations simultaneously.
Occupation records where the $a_t$ lie, but omits these products and
their shared coefficients. Fixed-factor solving, contractions, lift
conditions, and relaxations examined different parts of this gap.

The gap also appeared in verification. An early cardinality encoding
reused auxiliary variables between constraints and produced a
spurious contradiction. Qiushi identified the collision, withdrew the
affected exclusion, and rebuilt the encoding with disjoint counters.
The mathematical requirement was then explicit: every legal
multiplicity vector must extend to a simultaneous Boolean assignment.
This requirement survives in the formal copy and counter semantics
described in Appendix~\ref{app:certificate-realizations}.

The quotient search consequently left two distinct assets: the
restoration formula for upper bounds and valid occupation constraints
for lower bounds. The latter did not have to solve completion directly.
They could instead force a configuration in which the full tensor
itself supplied additional equations.

\subsection{Returning to the full tensor}
The split flattening introduced information absent from a quotient
support: each complete algorithm term contributes an image space of
dimension $\rank A_t$, and these spaces together must span dimension
$27$. After the geometric reduction, the profile $(16,1,3)$ used
exactly this dimension. There was no remaining dimensional overlap.
Qiushi then used the uniqueness of components in that direct sum to
obtain the saturation identities, and expanded their products in the
coordinates of matrix multiplication.

This was a change in the kind of conclusion being sought. Earlier
finite work had considered excluding completions of the surviving
profile by additional support systems. The saturation argument instead
gave a statement about every complete decomposition at rank sum $27$:
at most one first factor is invertible. It replaced the terminal
support exclusions by Proposition~\ref{prop:invertibility}. The
research note on the shortened proof also identified that the proposed
strengthening of orbit $479$ was unnecessary for the coset reduction:
its high-rank pair already has rank-one difference. That step requires
only the eight all-high plane raises. The existing bound
$R_{\F}(T_{W_{479}})\geq17$ still enters the line
cover~\eqref{eq:line-cover}.

The same statement then returned information to the constructive
problem. It restricts saturated $22$-term algorithms through
Corollary~\ref{cor:profiles}, while positive rank excess identifies
where the argument must be extended.

\subsection{Research memory and formal development}
Qiushi Engine's Meta-Trace links the history of questions, the evidence
for claims, and changes in the research trajectory~\cite{qiushi-optics}.
For this study, the accompanying record contains the deletion
diagnostic, the core restoration map,
the distinction between occupation and completion, the counter
correction, and the saturation argument, together with the notes,
programs, and computed data supporting them. Superseded interpretations
are retained alongside the evidence that corrected them.

The subsequent Lean development and manuscript preparation included
additional Codex assistance in proof completion, integration, and
verification. Formalization preserved the structural argument while
rebuilding its finite premises as proved quotient statements and
integer certificates. The library also includes the occupation,
substitution, and saturation lemmas used in the proof.
